% Letter to the children (two pages). No invented anecdote.
% TO BE PERSONALISED BY THE AUTHOR (à personnaliser par l'auteur): this letter is a frame; the author may
% add his own memories, which the agents are not allowed to invent.
\chapter*{Letter to my children}
\addcontentsline{toc}{chapter}{Letter to my children}
\markboth{Letter to my children}{Letter to my children}

\noindent\emph{Ilene, Alexis, Aurelia,}

\medskip
This book is about machines, galaxies and crystals, but it was written with you in mind. In a few lines, I would like
to tell you the three things this year taught me and that I would like to pass on to you.

\paragraph{Curiosity first} Everything told here began with a question: ``what if it were true?'', or ``what if it were
false?''. Asking questions is not the failing of a pupil who has not understood; it is the beginning of everything.
The machines in this book can answer almost anything. They do not yet know how to choose the questions that deserve to
be asked. That is your job, and it belongs to you. Keep the habit of asking ``why?'' and ``how do we know?'', even of
grown-ups, even of books, even of this one.

\paragraph{Checking your work} You have seen, in these pages, that a very learned machine can announce ``it's finished''
when nothing has been done. Humans can too, and I am the first. The only protection I know of is to check, and to have
someone else check. Redo a calculation another way. Try a case where you know the answer must be wrong, to see whether
you notice. Ask a friend to read it over. This is not mistrust: it is respect for what you do, and for the people who
are going to use it.

\paragraph{Owning up to your mistakes} This is the hardest lesson, and the most important. This year, I published things
that were wrong. I withdrew them, and I wrote down why, instead of quietly erasing them. It was not pleasant. But every
mistake I owned up to made the next piece of work better, and every hidden mistake would have made it more fragile. A
person who says ``I was wrong'' is not weaker than anyone else; they are more worthy of trust. You will have the right
to be wrong all your lives. What I wish for you is that you also have the courage to say so.

\medskip
There is something else I would like to tell you, something machines cannot do: a discovery only has meaning if it is
of use to someone. In this book, I have tried to choose questions that matter to people, and to say honestly what we
know and what we do not know. If one day you work with machines even more powerful than these, which is likely,
remember that it is up to you to decide what their work means.

At the very end of this book, five pictures are waiting for you. I will not tell you what they hide; you will find out.

\bigskip
\begin{flushright}
\emph{Papa}
\end{flushright}
